\chapter{Market making sur contrats à terme}\label{ch:hf:futures-market-making}

Dans un contrat à terme sur taux court, la file au meilleur cours peut contenir des dizaines de milliers de lots, et les exécutions sont partagées au prorata. Un teneur de marché qui
souhaite cent lots affiche davantage, car sa part de chaque exécution est sa taille rapportée à celle de tous, et chaque autre teneur raisonne de la même manière. Dans le
jeu de ce chapitre, dix teneurs qui souhaitent chacun 100 lots en affichent 230~; quinze en affichent 620~; vingt en affichent autant que la bourse le permet. Chacun est moins bien loti que si
tous avaient affiché ce qu'ils voulaient, et la profondeur au meilleur cours atteint 245 fois l'ordre moyen, comme Field et Large l'ont constaté sur de vrais
marchés au prorata.

\section{Carnets d'ordres au prorata et taille}

Le livre~1, chapitre~19 expose les règles d'allocation~: priorité temporelle (premier entré, premier exécuté), au prorata (chaque ordre au repos reçoit sa part de la
taille entrante), ainsi que les formules hybrides qui réservent d'abord une part à l'ordre de tête, c'est-à-dire au premier ordre à améliorer le prix, ou au teneur de marché principal. Sous priorité temporelle, la
décision consiste à choisir sa place dans la file (chapitre~5)~; au prorata, elle consiste à décider de la taille à afficher.

\begin{dated}{2026-09}{Comment une bourse répartit les exécutions}\label{dat:hf:futures-market-making:cme}
Selon le guide d'un fournisseur de données de marché décrivant en août 2025 les algorithmes d'appariement du CME, l'allocation au prorata apparie un agresseur entrant «~selon
le pourcentage au prorata de chaque ordre au repos~»~; les allocations sont arrondies à l'entier inférieur, zéro compris, et les lots en excédent sont
alloués selon la règle premier entré, premier sorti. Dans les algorithmes qui l'incluent, le premier ordre qui améliore le marché (l'ordre de tête) est apparié en premier, suivi de toute
allocation au teneur de marché principal. Le guide mentionnait le prorata pour certains carnets de calendar spread et un algorithme paramétrable pour les outrights, comme le
contrat à terme sur le Treasury note à deux ans.
\end{dated}

L'allocation d'un agresseur de 250 lots face à quatre ordres au repos de 100 (l'ordre de tête), 400, 300 et 200 montre la différence~: la priorité temporelle
donne 100 et 150 aux deux premiers~; le prorata pur donne 25, 100, 75 et 50~; l'ordre de tête puis le prorata donne 100, 67, 50 et 33 (\texttt{firm.match}).

\begin{definition}[Surcotation]\label{def:hf:futures-market-making:over}
La \emph{surcotation}\index{surcotation} (over-quoting) consiste à afficher, à un prix donné, davantage de taille qu'on ne souhaite en négocier, dans un carnet qui répartit les exécutions en
proportion de la taille, afin de recevoir une part plus grande de chaque exécution~; elle est rentable pour chaque teneur pris isolément et contre-productive pour tous ensemble.
\end{definition}

Le jeu (\cref{lst:hf:futures-market-making:game})~: $n$ teneurs de marché sont en attente au meilleur cours~; un agresseur de $X$ lots arrive, suivant une loi lognormale de
médiane 200~; le teneur $i$, qui affiche $s_i$ sur un total $S$, reçoit $s_i\min(X,S)/S$. Il gagne un tick par lot et paie $\gamma/2=0.005$ par lot au carré
d'inventaire. Seul, ou premier dans une file à priorité temporelle, il afficherait 100 lots, la taille dont l'exécution marginale vaut son risque marginal. Face aux autres,
afficher davantage achète une part plus grande de chaque ordre~; l'équilibre symétrique se trouve par meilleures réponses.

\begin{omfigure}
\begin{tikzpicture}
\begin{axis}[width=11cm,height=5.2cm,xlabel={\small teneurs au meilleur cours},ylabel={\small taille affichée / taille souhaitée},ymode=log,log basis y=10,
  xmin=0,xmax=21,ymin=0.8,ymax=130,ytick={1,2,5,10,20,50},yticklabels={1,2,5,10,20,50},tick label style={font=\footnotesize},grid=major,
  grid style={gray!25},table/col sep=comma]
\addplot[omThm,thick,mark=*,mark size=1.4pt] table[x=n,y=factor]{figdata/hft/15-futures-market-making/overquote.csv};
\draw[gray,dashed] (axis cs:0,50) -- (axis cs:21,50);
\node[font=\scriptsize,anchor=south] at (axis cs:10,50) {taille maximale autorisée par la bourse};
\end{axis}
\end{tikzpicture}

\omcaption{Facteur de \omterm{def:hf:futures-market-making:over}{surcotation}, soit la taille d'équilibre affichée par chaque teneur divisée par les 100 lots qu'il afficherait sous priorité temporelle, en fonction du
nombre de teneurs au meilleur cours (échelle logarithmique)~; agresseurs suivant une loi lognormale de médiane 200 lots, un tick gagné par lot, une \omterm{def:hf:inventory-models:penalty}{pénalité d'inventaire} $\gamma=0.01$
(0.005 par lot au carré), tailles plafonnées à 5\,000. Données~: \texttt{hf\_futures.game}.}\label{fig:hf:futures-market-making:over}
\end{omfigure}

Le facteur passe de 1.1 avec deux teneurs à 2.3 avec dix et 6.2 avec quinze~; avec vingt, la spirale ne s'arrête qu'au plafond
(\cref{fig:hf:futures-market-making:over}). L'utilité espérée de chaque teneur est plus faible que si tous avaient affiché 100~: 19.5 contre 22.3 avec dix teneurs.
Field et Large ont observé exactement ces symptômes sur quatre marchés de contrats à terme au prorata à un tick~: une profondeur aux cotations supérieure de deux
ordres de grandeur à l'ordre au marché moyen, et des taux d'annulation supérieurs à 96~\% des lots offerts. Ils ont modélisé la cause comme la complémentarité stratégique du jeu~: chaque
trader risque de trop négocier avec un ordre surdimensionné qu'il s'attend surtout à annuler.

Le risque tient à l'ordre rare et de grande taille. Un agresseur au 99.9\up{e}~centile des tailles simulées (7\,907 lots) remplit chacun des dix teneurs avec 230 lots,
soit 2.3 fois ce que chacun voulait, tous à la fois et tous du même côté~: l'ensemble du marché se retrouve long d'un seul coup.

\section{Calendar spreads et prix implicites}

Un calendar spread (le livre~1, chapitre~19) achète une échéance et en vend une autre. Les bourses cotent des carnets de spreads à côté des outrights et les relient par
prix implicites~: le prix implied-in d'un spread provient des deux carnets d'outrights (cours acheteur de l'échéance proche moins cours vendeur de l'échéance lointaine), le prix implied-out
d'un outright se déduit d'un carnet de spread et de l'autre outright. Le moteur d'appariement de la bourse peut traverser ces liens, si bien qu'un ordre dans un carnet peut être
exécuté contre des ordres dans deux autres.

\section{Arbitrage entre carnets de spreads et d'outrights}

\begin{definition}[Arbitrage de prix implicite]\label{def:hf:futures-market-making:implied}
L'\emph{arbitrage de prix implicite}\index{arbitrage de prix implicite} (implied-price arbitrage) consiste à négocier un carnet de spread contre ses deux carnets d'outrights lorsque le prix direct du spread
est en dehors de l'intervalle implicite des outrights (ou l'inverse), avant que l'appariement implicite de la bourse ou d'autres intervenants ne comblent l'écart.
\end{definition}

Là où la bourse calcule des prix implicites pour chaque lien, beaucoup de ces écarts sont comblés par le moteur lui-même~; là où elle ne le fait pas (pour certains spreads,
entre bourses, ou entre un spread et un pack de plusieurs outrights), ils sont comblés par les traders. Les écarts viennent de retards~: les teneurs de marché du carnet de spread
cotent autour de la différence entre les outrights, telle qu'ils l'ont vue en dernier. Dans le strip de quatre contrats à terme trimestriels de ce chapitre, les carnets d'outrights sont larges d'un tick et
les carnets de spreads, larges eux aussi d'un tick, autour d'une différence retardée~; des arbitrages apparaissent lorsque le retard couvre assez d'événements de données de marché pour que la pente de la courbe varie de
deux ticks~: aucun avec un retard d'un événement, 0.9 pour mille événements de spread avec cinq, 5.9 avec dix et 50 avec trente, chacun valant environ un tick
(\cref{fig:hf:futures-market-making:implied}).

\begin{omfigure}
\begin{tikzpicture}
\begin{axis}[width=11cm,height=4.8cm,xlabel={\small retard des carnets de spreads (événements de données de marché)},ylabel={\small arbitrages pour 1\,000},xmin=0,xmax=31,
  ymin=0,ymax=55,tick label style={font=\footnotesize},grid=major,grid style={gray!25},table/col sep=comma]
\addplot[omThm,thick,mark=*,mark size=1.4pt] table[x=lag,y=per1000]{figdata/hft/15-futures-market-making/implied.csv};
\end{axis}
\end{tikzpicture}

\omcaption{\omterm{def:hf:futures-market-making:implied}{Arbitrages de prix implicite} entre trois carnets de calendar spreads et quatre carnets d'outrights, pour mille observations de carnets de spreads, en fonction du
nombre d'événements de données de marché dont les cotations des carnets de spreads retardent sur les outrights~; le niveau et la pente de la courbe évoluent de 0.6 et 0.4 tick par événement~; 5\,000 événements. Données~:
\texttt{hf\_futures.implied}.}\label{fig:hf:futures-market-making:implied}
\end{omfigure}

\section{Strips de contrats à terme sur taux}

Un strip de contrats à terme sur taux court (le livre~2, chapitre~8) est une courbe de contrats trimestriels, chacun étant un taux forward pour une période de trois mois, négocié
comme outrights, calendar spreads, butterflies, packs et bundles. Un teneur de marché du strip les cote tous à partir d'un même modèle de courbe (quelques
facteurs~: niveau, pente, courbure, plus les dates de réunion de la banque centrale) et couvre toute exécution avec l'instrument le moins cher. Les carnets du strip
sont au prorata ou hybrides sur plusieurs marchés, si bien que la \omterm{def:hf:futures-market-making:over}{surcotation} décrite dans les sections précédentes s'applique à chacun~; le modèle de la courbe est ce qui garde
les cotations du teneur sur des centaines de carnets cohérentes entre elles.

\section{Spreads inter-commodités}

Deux contrats liés (deux points d'une courbe de rendement sur des contrats différents, le brut et ses produits, le maïs et l'éthanol) se négocient comme un
spread inter-commodités lorsque leurs prix partagent un facteur. Le crédit de marge accordé par la bourse pour le spread (le livre~1, chapitre~20) abaisse son coût de
portage. Un teneur de marché cote le spread à partir de son modèle de la relation et couvre chaque exécution dans la jambe la moins chère à négocier~; le risque est la
dérive propre à la relation, que nulle couverture ne supprime.

\section{Fiches de stratégie}

\begin{strategyfile}{Cotation outright au prorata}\label{strat:hf:futures-market-making:prorata}
\sfield{Qui paie le teneur de marché, et pourquoi} Hedgers et traders directionnels qui traversent une fourchette d'un tick en taille.
\sfield{Instruments et plateformes} Outrights sur contrats à terme sur taux court et sur obligations, sur carnets au prorata ou hybrides.
\sfield{Signal} La fair value à l'intérieur du tick (chapitre~2) et la probabilité d'un gros agresseur, déduites du flux d'ordres et des autres carnets.
\sfield{Taille et exécution} Afficher la taille d'équilibre, et non la taille souhaitée~; annuler avant les gros ordres prévus~; viser la priorité au premier ordre et celle du teneur de
marché principal, là où elle est disponible.
\sfield{Coûts} Frais de la bourse par lot, inventaire issu des grosses exécutions, limites de messages dues à l'annulation de la plus grande partie de ce qui est affiché.
\sfield{Mode d'échec} Le très gros ordre, rare~: dix teneurs de marché exécutés chacun à 2.3 fois leur taille souhaitée par un seul agresseur.
\sfield{Horizon, capacité, infrastructure} De la seconde à la minute~; capacité importante~; un modèle de courbe et une annulation rapide.
\sfield{Backtest en toute honnêteté} L'algorithme d'allocation et ses paramètres par produit~; les tailles des autres (non observables) modélisées, et non supposées sans effet.
\sfield{Sources} Field et Large (2008)~; l'encadré daté.
\end{strategyfile}

\begin{strategyfile}{Cotation de calendar spread avec jambes implicites}\label{strat:hf:futures-market-making:spread}
\sfield{Qui paie le teneur de marché, et pourquoi} Les hedgers qui font le roll de leurs positions d'une échéance à la suivante, et qui négocient le spread pour éviter le \omterm{def:hf:index-and-basket-arbitrage:legging}{risque d'exécution par jambes}.
\sfield{Instruments et plateformes} Les carnets de calendar spreads de la bourse et leurs outrights.
\sfield{Signal} L'écart des outrights au modèle de courbe, les prix implicites de la bourse servant de cotation concurrente.
\sfield{Taille et exécution} Coter le spread à l'intérieur du marché implicite~; couvrir une exécution dans les outrights lorsque le modèle indique que le spread ne
reviendra pas rapidement vers sa moyenne.
\sfield{Coûts} Frais sur le spread et sur toute couverture~; position dans la file derrière les ordres implicites.
\sfield{Mode d'échec} Les périodes de roll concentrent flux et concurrents~; le moteur implicite devance une cotation lente.
\sfield{Horizon, capacité, infrastructure} De la seconde au jour~; calendriers de roll.
\sfield{Backtest en toute honnêteté} Les prix implicites que le moteur aurait générés~; les volumes des périodes de roll.
\sfield{Sources} Le livre~1, chapitre~19~; le présent chapitre.
\end{strategyfile}

\begin{strategyfile}{Arbitrage de prix implicite entre carnets de spreads et carnets d'outrights}\label{strat:hf:futures-market-making:implied}
\sfield{Qui paie le teneur de marché, et pourquoi} Les teneurs de marché des carnets de spreads dont les prix sont en retard sur les outrights.
\sfield{Instruments et plateformes} Un carnet de spreads et ses deux carnets d'outrights, où l'appariement implicite de la bourse ne ferme pas tous les liens.
\sfield{Signal} Le prix direct du spread hors de l'intervalle implicite issu des outrights.
\sfield{Taille et exécution} Prendre le spread et les deux jambes outrights en même temps~; dimensionner à la plus petite des trois quantités au meilleur cours.
\sfield{Coûts} Frais sur les trois carnets~; \omterm{def:hf:index-and-basket-arbitrage:legging}{risque d'exécution par jambes} si une jambe est manquée (chapitre~13).
\sfield{Mode d'échec} Le moteur implicite de la bourse et les traders plus rapides~: dans la simulation, les écarts n'apparaissent que lorsque les carnets de spreads accusent un retard de plusieurs
événements.
\sfield{Horizon, capacité, infrastructure} Microsecondes~; les données des trois carnets sur une même horloge.
\sfield{Backtest en toute honnêteté} Événements de carnet séquencés sur les trois carnets~; savoir si le moteur implicite aurait apparié en premier.
\sfield{Sources} Le présent chapitre.
\end{strategyfile}

\begin{strategyfile}{Cotation d'un strip de contrats à terme sur taux court}\label{strat:hf:futures-market-making:strip}
\sfield{Qui paie le teneur de marché, et pourquoi} Hedgers de taux et spéculateurs négociant n'importe quel point du strip, ainsi que les packs et les bundles.
\sfield{Instruments et plateformes} Un strip de contrats à terme trimestriels sur taux court, ses spreads, ses butterflies, ses packs et ses bundles (le livre~2, chapitre~8).
\sfield{Signal} Un modèle de courbe~: niveau, pente, courbure, dates des réunions de la banque centrale, calibré sur l'ensemble des carnets du strip à la fois.
\sfield{Taille et exécution} Coter chaque carnet à partir du modèle~; couvrir toute exécution dans l'instrument le moins cher qui supprime son exposition à la courbe.
\sfield{Coûts} Frais, risque d'inventaire de la \omterm{def:hf:futures-market-making:over}{surcotation} dans les carnets au prorata, erreur de modèle autour des réunions de la banque centrale.
\sfield{Mode d'échec} Les surprises des banques centrales déplacent la courbe de façons que les facteurs du modèle ne couvrent pas.
\sfield{Horizon, capacité, infrastructure} De la seconde au jour~; des centaines de carnets issus d'un même modèle.
\sfield{Backtest en toute honnêteté} Le calendrier des réunions tel qu'il était connu à l'époque~; les règles d'allocation de tous les carnets.
\sfield{Sources} Le livre~2, chapitre~8.
\end{strategyfile}

\begin{strategyfile}{Négociation de spreads inter-commodités}\label{strat:hf:futures-market-making:intercommodity}
\sfield{Qui paie le teneur de marché, et pourquoi} Hedgers de marges de transformation et de positions sur la courbe, qui négocient le spread comme un seul instrument.
\sfield{Instruments et plateformes} Deux contrats à terme liés et, lorsqu'il est coté, leur carnet de spreads inter-commodités.
\sfield{Signal} L'écart du spread à un modèle de la relation entre les deux prix.
\sfield{Taille et exécution} Coter le spread~; couvrir les exécutions dans la jambe offrant la meilleure liquidité.
\sfield{Coûts} Frais des deux jambes, marge (réduite par le crédit de marge accordé par la bourse pour le spread), dérive de la relation.
\sfield{Mode d'échec} Changement structurel de la relation (un arrêt de raffinerie, une nouvelle règle de livraison).
\sfield{Horizon, capacité, infrastructure} De la minute à la semaine.
\sfield{Backtest en toute honnêteté} Crédits de marge tels qu'ils étaient~; changements de livraison et de contrat à leur date.
\sfield{Sources} Le livre~1, chapitre~20.
\end{strategyfile}

\section{Tutoriel~: afficher dix fois ce que l'on veut}\label{tut:hf:futures-market-making:tut}

\begin{tutorial}
\textbf{Objectif.} Trouver la taille d'équilibre qu'affichent les teneurs de marché dans un carnet au prorata, son coût pour eux, et les arbitrages implicites qu'un carnet de spreads en retard laisse subsister.
\textbf{État final~:} les \cref{fig:hf:futures-market-making:over,fig:hf:futures-market-making:implied}.
\begin{tutsteps}
\item \textbf{Allocation} d'un agresseur selon trois règles avec \texttt{firm.match} (\texttt{hf\_futures.allocation\_example}).
\item \textbf{Le jeu}~: meilleures réponses à la taille totale des autres, amorties, jusqu'à stabilisation.
\omcode{code/firm/futmm/firm_futmm.py}{45}{62}{Taille optimale d'un teneur de marché face au total des autres~; équilibre symétrique obtenu par meilleures réponses amorties.}\label{lst:hf:futures-market-making:game}
\item \textbf{Stress}~: les exécutions qu'un agresseur au 99.9\up{e}~centile donne à chaque teneur de marché.
\item \textbf{Le strip}~: quatre contrats à terme trimestriels, trois carnets de spreads cotant une différence avec retard~; balayage à la recherche d'arbitrages implicites.
\omcode{code/firm/futmm/firm_futmm.py}{99}{114}{Cotations directes du spread comparées au spread implicite issu des outrights, avec \texttt{firm.match.implied\_in}.}\label{lst:hf:futures-market-making:scan}
\end{tutsteps}
\textbf{Pour aller plus loin.} Ajouter une allocation au premier ordre de 40~\% au jeu et voir qui en bénéficie~; laisser un teneur de marché prévoir les gros agresseurs et annuler
avant eux~; faire tourner le carnet au prorata de \texttt{firm.exchsim} sur son préréglage \texttt{pro\_rata\_futures}.
\end{tutorial}

\begin{center}
\small
\begin{tabular}{@{}lrrrrrr@{}}
\toprule
teneurs de marché & 2 & 5 & 10 & 12 & 15 & 20\\
\midrule
taille affichée (lots) & 110 & 143 & 230 & 298 & 619 & 5\,000\\
facteur de \omterm{def:hf:futures-market-making:over}{surcotation} & 1.1 & 1.4 & 2.3 & 3.0 & 6.2 & 50\\
utilité, équilibre & 41.2 & 29.6 & 19.5 & 16.9 & 13.9 & 12.3\\
utilité, tous affichant 100 & 41.4 & 31.2 & 22.3 & 20.0 & 17.4 & 14.3\\
exécution d'un ordre de 7\,907 lots & 110 & 143 & 230 & 298 & 527 & 395\\
profondeur / ordre moyen & 0.5 & 1.8 & 5.6 & 8.8 & 23 & 245\\
\bottomrule
\end{tabular}
\end{center}

\section{Construction~: le module de market making sur contrats à terme}\label{bld:hf:futures-market-making:futmm}

\begin{build}
\bfield{Rôle} Modéliser les exécutions au prorata et l'équilibre de \omterm{def:hf:futures-market-making:over}{surcotation}, et balayer un strip à la recherche d'\omterm{def:hf:futures-market-making:implied}{arbitrages de prix implicite}.
\bfield{Interface} \texttt{fills(s, X)}, \texttt{utility}, \texttt{best\_response}, \texttt{equilibrium(n, X, h, gamma, grid)},
\texttt{fifo\_size}, \texttt{aggressors}, \texttt{Strip(n\_contracts, steps, seed, lag)} avec \texttt{outright(t, i)} et \texttt{spread(t, i)},
\texttt{scan(strip)}. Construit sur \texttt{firm.match} (le livre~1)~: \texttt{pro\_rata}, \texttt{configurable}, \texttt{implied\_in}.
\bfield{Règles} Les exécutions sont des parts continues de $\min(X,S)$~; le haut de la grille est la taille maximale d'ordre de la bourse.
\bfield{Tests d'acceptation} \texttt{code/firm/futmm/tests/}~: exécutions continues à l'arrondi près de \texttt{firm.match.pro\_rata}~; sans concurrence, la meilleure réponse
est la taille de priorité temporelle~; la meilleure réponse croît avec la foule~; l'équilibre croît avec le nombre de teneurs et les laisse moins bien lotis
que s'ils affichaient ce qu'ils veulent~; pas d'arbitrage sans retard, davantage avec plus de retard.
\bfield{Pour aller plus loin} Tranches au premier ordre et du teneur de marché principal dans le jeu~; le jeu sur le préréglage au prorata de \texttt{firm.exchsim}~; packs et butterflies dans
le scanner.
\end{build}

\begin{omsources}
\item J.~Field, J.~Large, Pro-rata matching and one-tick futures markets, CFS Working Paper 2008/40, 2008.
\item Databento, CME matching algorithms explained, 1\up{er} août 2025.
\end{omsources}

\section{Exercices}

\begin{exercise}[$\star$]\label{exo:hf:futures-market-making:1}
Quatre ordres de 100 (le premier ordre), 400, 300 et 200 lots reposent à un même prix. Allouer un agresseur de 250 selon la priorité temporelle, le prorata pur, et l'allocation au
premier ordre puis au prorata.
\end{exercise}

\begin{exercise}[$\star$]\label{exo:hf:futures-market-making:2}
Dix teneurs de marché affichent chacun 230 lots. Que donne un agresseur de 1\,000 lots à chacun, et un de 7\,907~?
\end{exercise}

\begin{exercise}[$\star$]\label{exo:hf:futures-market-making:3}
Le contrat de tête a un cours acheteur de 9\,612 et un cours vendeur de 9\,613, le suivant un cours acheteur de 9\,598 et un cours vendeur de 9\,599. Quel est le marché implicite du spread~? Un carnet de spread est offert à
12~: que faire~?
\end{exercise}

\begin{exercise}[$\star\star$]\label{exo:hf:futures-market-making:4}
Expliquer pourquoi l'équilibre est un dilemme du prisonnier.
\end{exercise}

\begin{exercise}[$\star\star$]\label{exo:hf:futures-market-making:5}
Expliquer pourquoi les carnets au prorata ont des taux d'annulation supérieurs à 96~\%.
\end{exercise}

\begin{exercise}[$\star\star$]\label{exo:hf:futures-market-making:6}
Expliquer comment l'allocation au premier ordre modifie les incitations d'un teneur de marché.
\end{exercise}

\begin{exercise}[$\star\star\star$]\label{exo:hf:futures-market-making:7}
\emph{Programmation.} Doubler la \omterm{def:hf:inventory-models:penalty}{pénalité d'inventaire} (\texttt{GAMMA = 0.02}). Quelle taille afficherait un teneur de marché sous la priorité temporelle, et qu'affichent dix teneurs de marché~?
\end{exercise}

\begin{exercise}[$\star\star\star$]\label{exo:hf:futures-market-making:8}
\emph{Trouver la faille.} «~La profondeur au meilleur niveau est de 20\,000 lots, donc notre ordre de 500 lots ne fera pas bouger le prix et nous obtiendrons 500 lots.~»
\end{exercise}

\section{Problème~: afficher dix fois ce que l'on veut}

\begin{problem}[{Problème du week-end --- afficher dix fois ce que l'on veut}]\label{pb:hf:futures-market-making:1}
Les teneurs de marché d'un carnet de contrats à terme sur taux courts au prorata choisissent combien afficher, et un carnet de spread voisin accuse un retard.

\medskip\noindent\textbf{Partie I --- Allocation.}
\begin{enumerate}
\item Décrire la priorité temporelle, le prorata et l'allocation au premier ordre, ainsi que la description d'une bourse donnée dans l'encadré daté.
\item Allouer l'agresseur de l'exemple selon chaque règle.
\item Définir la \omterm{def:hf:futures-market-making:over}{surcotation}.
\item Qu'ont observé Field et Large sur les marchés au prorata~?
\end{enumerate}

\medskip\noindent\textbf{Partie II --- Le jeu.}
\begin{enumerate}[resume]
\item Écrire l'exécution et l'utilité d'un teneur de marché.
\item Quelle taille un teneur de marché affiche-t-il sous la priorité temporelle, et pourquoi~?
\item Comment détermine-t-on l'équilibre symétrique~?
\item Pourquoi la spirale ne s'arrête-t-elle qu'au plafond avec vingt teneurs de marché~?
\end{enumerate}

\medskip\noindent\textbf{Partie III --- Mesures.}
\begin{enumerate}[resume]
\item Donner le facteur de \omterm{def:hf:futures-market-making:over}{surcotation} pour 2, 10, 15 et 20 teneurs de marché.
\item Comparer l'utilité à l'équilibre avec celle obtenue lorsque tous affichent 100.
\item Que fait l'agresseur au 99.9\up{e}~centile à chaque teneur de marché~?
\item Définir l'\omterm{def:hf:futures-market-making:implied}{arbitrage de prix implicite} et donner les décomptes selon le retard.
\end{enumerate}

\medskip\noindent\textbf{Partie IV --- Le verdict.}
\begin{enumerate}[resume]
\item Énoncer le \emph{résultat nommé}~: le facteur de \omterm{def:hf:futures-market-making:over}{surcotation} à l'équilibre dans un carnet au prorata, et le risque d'inventaire qu'il crée lorsqu'un ordre
d'une taille inattendue exécute tous les teneurs à la fois.
\item Comment un teneur de marché doit-il dimensionner son affichage, sachant l'équilibre~?
\item Que pourrait modifier une bourse pour arrêter la spirale~?
\item Pourquoi les teneurs de marché de strip cotent-ils des centaines de carnets à partir d'un même modèle~?
\item Quand un spread entre matières premières cesse-t-il d'être une couverture~?
\item Laquelle des cinq fiches de stratégie dépend le plus de la vitesse~?
\item Comment réaliser le backtest d'une stratégie au prorata lorsque les tailles des autres ne sont pas observables~?
\item En une phrase~: que rapporte la taille dans un carnet au prorata~?
\end{enumerate}
\end{problem}

\section{Questions d'entretien}

\begin{interviewq}[$\star$ \iqroles{trader}]\label{iq:hf:futures-market-making:1}
Pourquoi la fourchette d'un contrat à terme sur taux court vaut-elle presque toujours un tick~?
\end{interviewq}

\begin{interviewq}[$\star\star$ \iqroles{chercheur}]\label{iq:hf:futures-market-making:2}
Dérivez la meilleure réponse dans le jeu au prorata lorsque les agresseurs n'épuisent jamais le carnet.
\end{interviewq}

\begin{interviewq}[$\star\star$ \iqroles{trader}]\label{iq:hf:futures-market-making:3}
Vous affichez 2\,000 lots et un ordre de 10\,000 lots frappe le cours acheteur. Que se passe-t-il ensuite, et que doit faire votre système~?
\end{interviewq}

\begin{interviewq}[$\star\star$ \iqroles{développeur}]\label{iq:hf:futures-market-making:4}
Implémentez l'allocation au prorata avec arrondi à l'inférieur, une allocation minimale et les restes en FIFO. Quels sont les cas limites~?
\end{interviewq}

\begin{interviewq}[$\star\star$ \iqroles{risque}]\label{iq:hf:futures-market-making:5}
La taille affichée par votre société dans un carnet au prorata est cinquante fois sa limite de position. Est-ce une violation~?
\end{interviewq}

\begin{interviewq}[$\star\star\star$ \iqroles{chercheur}]\label{iq:hf:futures-market-making:6}
Comment estimeriez-vous, à partir des données publiques de marché, la part de la profondeur affichée dans un carnet au prorata qui relève de la \omterm{def:hf:futures-market-making:over}{surcotation}~?
\end{interviewq}
